\section{Preliminaries and Problem Formulation}
\label{sec:preliminaries}

\noindent\textbf{Tabular In-Context Prediction.}\label{sec:tabular_icl}
Let $\mathcal{D}_s=\{(x_j^s,y_j^s)\}_{j=1}^{n_s}$ and
$\mathcal{X}_t=\{x_i^t\}_{i=1}^{n_t}$ denote labeled reference data and
unlabeled target queries. Indices $j$ and $i$ refer to reference rows and
target queries, respectively, with $x_j^s,x_i^t\in\mathcal{X}\subseteq
\mathbb{R}^d$. We take $\mathcal{Y}=\{1,\ldots,L\}$ for classification and
$\mathcal{Y}\subseteq\mathbb{R}$ for regression; target labels are unavailable
at inference and used only for evaluation.
A tabular in-context learner predicts each query by conditioning on labeled
rows supplied as context. For $x_i^t$, let
$C_i=(X_i^{\mathrm{ctx}},\mathbf{y}_i^{\mathrm{ctx}})
\in\mathbb{R}^{B_i\times d}\times\mathcal{Y}^{B_i}$ be its context of $B_i$
labeled rows. Thus, $\mathcal{D}_s$ is the available reference set, whereas
$C_i$ is the model input selected or constructed from the available data.
Let $f_\theta$ denote a TFM pretrained before deployment via synthetic prior
fitting or real-table pretraining
\citep{hollmann2025tabpfn,qu2025tabicl,ma2025tabdpt} and frozen at inference.
Given $C_i$ and $x_i^t$, it returns an $L$-class probability vector,
\begin{equation}
  p_\theta(\cdot\mid x_i^t,C_i)=f_\theta(C_i,x_i^t).
  \label{eq:tfm_inference}
\end{equation}
For regression, $f_\theta$ instead returns a real-valued prediction or
predictive distribution. Section~\ref{sec:method} extends this interface to
context ensembles.

\noindent\textbf{Prediction under Data Change.}\label{sec:data_change}
Let $P_s(X,Y)$ and $P_t(X,Y)$ denote the joint distributions of the reference
and target environments, respectively, with
$(x_j^s,y_j^s)\sim P_s$ and $(x_i^t,y_i^t)\sim P_t$. We call $P_t=P_s$ the
nominal setting. A \emph{controlled source-context perturbation} intentionally
modifies the accessible labeled evidence to isolate a particular Context
failure, without requiring that the operation correspond to a named deployment
shift. A \emph{distribution shift} instead describes a difference between the
underlying source and target distributions, $P_t\neq P_s$. Controlled
perturbations modify the source context. Distribution-shift operators may act
on the source, the target, or both, according to the family definition in
Appendix~\ref{app:data_transformations}.
For a predictor $g$ and loss $\ell$, performance in the target environment is
measured by the target risk
\begin{equation}
  \mathcal{R}_t(g)
  =\mathbb{E}_{(X,Y)\sim P_t}\!\left[\ell\!\left(g(X),Y\right)\right].
  \label{eq:target_risk}
\end{equation}
Our objective is to reduce target risk across nominal and changed conditions
using $\mathcal D_s$ and the unlabeled queries $\mathcal X_t$, while keeping
$f_\theta$ fixed.
